%%====================================================================
%% OFI_QM_Foundations.tex
%%
%% Ordered Fractional Integration and the Foundations of Quantum Mechanics:
%%   Uncertainty, Entanglement, Decoherence, and Measurement
%%
%% Author: Joseph R. Escamilla
%% Atlas SDK · Senpai Productions · Mathematical Sciences Division
%% Date: April 2026
%%
%% Cross-references:
%%   OFI.tex                       -- core Riesz potential / semigroup framework
%%   OFI_Natural_Derivation.tex    -- kinetic-symbol principle alpha = k/2
%%   OFI_Cosmological_Constant.tex -- OFI vacuum arms (-1, +1/12)
%%   OFI_Cramer_TI_Tachyon.tex     -- tachyon arm, imaginary defect, transactions
%%   OFI_ADM_WheelerDeWitt.tex     -- OFI Hilbert space, essential self-adjointness
%%   OFI_Unruh_Effect.tex          -- null-hinge transparency, Bogoliubov structure
%%
%% Epistemic convention (used throughout):
%%   \DERIVED{} -- logically derived from OFI axioms + standard QM/QFT
%%   \INPUT{}   -- requires experimental or prior-paper input
%%   \OPEN{}    -- genuine open problem; not claimed
%%====================================================================

\documentclass[12pt]{article}

\usepackage{amsmath,amssymb,amsthm}
\usepackage[margin=1.25in]{geometry}
\usepackage{hyperref}
\usepackage{xcolor}
\usepackage{booktabs}
\usepackage{array}
\usepackage{microtype}
\usepackage{lmodern}
\usepackage{mathrsfs}
\usepackage{enumitem}
\usepackage{tcolorbox}
\tcbuselibrary{skins,breakable}

%%--------------------------------------------------------------------
%% Hyperref
%%--------------------------------------------------------------------
\hypersetup{
  pdftitle  = {OFI and the Foundations of Quantum Mechanics},
  pdfauthor = {Joseph R.\ Escamilla},
  colorlinks= true,
  linkcolor = blue!60!black,
  citecolor = green!50!black,
  urlcolor  = blue!70!black
}

%%--------------------------------------------------------------------
%% Theorem environments
%%--------------------------------------------------------------------
\newtheorem{theorem}{Theorem}[section]
\newtheorem{lemma}[theorem]{Lemma}
\newtheorem{proposition}[theorem]{Proposition}
\newtheorem{corollary}[theorem]{Corollary}
\newtheorem{definition}[theorem]{Definition}
\newtheorem{remark}[theorem]{Remark}
\theoremstyle{definition}
\newtheorem{claim}[theorem]{Claim}
\newtheorem{conjecture}[theorem]{Conjecture}
\newtheorem{example}[theorem]{Example}

%%--------------------------------------------------------------------
%% Notation
%%--------------------------------------------------------------------
\newcommand{\R}{\mathbb{R}}
\newcommand{\C}{\mathbb{C}}
\newcommand{\N}{\mathbb{N}}
\newcommand{\abs}[1]{\left|#1\right|}
\newcommand{\norm}[1]{\left\|#1\right\|}
\newcommand{\inner}[2]{\langle #1,\, #2 \rangle}
\newcommand{\bra}[1]{\langle #1 \rvert}
\newcommand{\ket}[1]{\lvert #1 \rangle}
\newcommand{\braket}[2]{\langle #1 \mid #2 \rangle}
\newcommand{\expect}[1]{\langle #1 \rangle}
\newcommand{\OFI}{\mathrm{OFI}}
\newcommand{\Ia}{I^{\alpha}}
\newcommand{\Ihalf}{I^{1/2}}
\newcommand{\Ione}{I^{1}}
\newcommand{\Id}{\mathrm{Id}}
\newcommand{\Tr}{\operatorname{Tr}}
\newcommand{\sgn}{\operatorname{sgn}}

%% OFI defect
\newcommand{\Rdef}[1]{R^{#1}}
\newcommand{\RdefA}[2]{R^{#1}_{#2}}

%%--------------------------------------------------------------------
%% Epistemic flags
%%--------------------------------------------------------------------
\definecolor{derivedcolor}{RGB}{0,110,0}
\definecolor{inputcolor}{RGB}{160,0,0}
\definecolor{opencolor}{RGB}{0,0,180}
\newcommand{\DERIVED}[1]{\textcolor{derivedcolor}{\textbf{[DERIVED: #1]}}}
\newcommand{\INPUT}[1]{\textcolor{inputcolor}{\textbf{[INPUT: #1]}}}
\newcommand{\OPEN}[1]{\textcolor{opencolor}{\textbf{[OPEN: #1]}}}

%%--------------------------------------------------------------------
%% Coloured summary boxes
%%--------------------------------------------------------------------
\tcbset{
  resultbox/.style={
    colback=green!5!white, colframe=green!50!black,
    fonttitle=\bfseries, breakable, enhanced
  },
  openbox/.style={
    colback=blue!5!white, colframe=blue!50!black,
    fonttitle=\bfseries, breakable, enhanced
  },
  warningbox/.style={
    colback=red!5!white, colframe=red!60!black,
    fonttitle=\bfseries, breakable, enhanced
  }
}

%%====================================================================
\begin{document}
%%====================================================================

\title{%
  \textbf{Ordered Fractional Integration and the Foundations%
          of Quantum Mechanics}\\[0.5em]
  \large Uncertainty, Entanglement, Decoherence, and Measurement%
         from the OFI Defect Structure%
}

\author{%
  Joseph R.\ Escamilla\\[0.3em]
  \small Atlas SDK $\cdot$ Senpai Productions\\
  \small Mathematical Sciences Division
}

\date{April 2026}

\maketitle

\begin{abstract}
We examine four foundational pillars of quantum mechanics through the lens
of Ordered Fractional Integration (OFI), realized as the Riesz potential
$I^\alpha$ with Fourier multiplier $|\xi|^{-\alpha}$ and semigroup law
$I^\alpha I^\beta = I^{\alpha+\beta}$.  The OFI defect $R^\alpha[\varphi]
= \mathcal{L}[\varphi] - \mathcal{L}[I^\alpha\varphi]$ measures the
failure of a Lagrangian to commute with $I^\alpha$.

The established OFI commutator $[I\circ D,\, D\circ I] = \tfrac{1}{12}$
(a structural consequence of the Euler--Maclaurin formula and Bernoulli
regularization) is the central object.

\textbf{Section~\ref{sec:uncertainty}} (Uncertainty Principle): We derive a
Robertson-type uncertainty inequality $\Delta A\,\Delta B \geq \tfrac{1}{24}$
directly from the OFI commutator, and we carry out a careful dimensional
analysis of the relationship between the OFI quantum $\tfrac{1}{12}$ and
Planck's constant $\hbar$.  We are explicit about what is and is not derived.

\textbf{Section~\ref{sec:entanglement}} (Bipartite Entanglement): We define the
OFI entanglement defect $\Delta R$ for bipartite systems and show $\Delta R = 0$
iff the state is a product state.  We present the honest status of OFI vis-\`{a}-vis
the CHSH Bell bound $2\sqrt{2}$.

\textbf{Section~\ref{sec:decoherence}} (Decoherence): We show that diagonal
density-matrix elements have real OFI defect (stable) while off-diagonal elements
acquire imaginary defect (oscillatory, $\to 0$ under RG flow from the tachyon arm).
This identifies a natural pointer-state basis in the OFI framework.

\textbf{Section~\ref{sec:measurement}} (Measurement and Collapse): We formulate
the OFI hypothesis for wavefunction collapse as an RG flow from the tachyon arm
($\alpha = -1$) to the physical vacuum ($\alpha = +\tfrac{1}{12}$), and we
derive the Born rule from OFI defect norms combined with the decoherence result
of Section~\ref{sec:decoherence}.

Throughout, \DERIVED{} marks results that follow from OFI axioms, \INPUT{} marks
results that require external input, and \OPEN{} marks genuine open problems.
\end{abstract}

\tableofcontents
\newpage

%%====================================================================
\section{Preliminaries: OFI Framework}
\label{sec:prelim}
%%====================================================================

\subsection{Riesz Potential as OFI Operator}

\begin{definition}[OFI / Riesz Potential]
\label{def:riesz}
For $n \geq 1$ and $0 < \alpha < n$, the \emph{Ordered Fractional
Integration} operator at order $\alpha$ in $\R^n$ is the Riesz potential
\begin{equation}
  (I^\alpha f)(x)
  \;=\;
  c_{n,\alpha}
  \int_{\R^n} |x - y|^{\alpha - n}\, f(y)\,dy,
  \qquad
  c_{n,\alpha}
  = \frac{\Gamma\!\bigl(\tfrac{n-\alpha}{2}\bigr)}
         {2^{\alpha}\,\pi^{n/2}\,\Gamma\!\bigl(\tfrac{\alpha}{2}\bigr)}.
\end{equation}
Key properties:
\begin{itemize}[leftmargin=2em]
\item \textbf{Fourier multiplier:}
  $\widehat{I^\alpha f}(\xi) = |\xi|^{-\alpha}\,\hat{f}(\xi)$.
\item \textbf{Semigroup:}
  $I^\alpha I^\beta = I^{\alpha+\beta}$ for
  $\alpha,\beta,\alpha+\beta \in (0,n)$.
\item \textbf{Identity at $\alpha = 0$:}
  $I^0 = \Id$ (by analytic continuation; the $\Gamma(0)$ pole cancels
  the kernel singularity to produce $\delta(x-y)$).
\item \textbf{Symmetry:}
  $|x-y| = |y-x|$; no preferred boundary or left endpoint.
\end{itemize}
\end{definition}

\begin{definition}[OFI Defect]
\label{def:defect}
Let $\mathcal{L}[\varphi]$ be a Lagrangian or functional, and let
$\varphi_a := I^\alpha\varphi$ be the OFI-filtered field.  The
\emph{OFI defect} at order $\alpha$ is
\begin{equation}
  R^\alpha[\varphi]
  \;:=\;
  \mathcal{L}[\varphi] - \mathcal{L}[\varphi_a].
  \label{eq:defect}
\end{equation}
$R^\alpha[\varphi]$ measures how much $\mathcal{L}$ fails to commute with
$I^\alpha$.
\end{definition}

\subsection{The OFI Commutator}

Let $I = I^1$ denote integration (order 1 in 1D) and $D = d/dx$.
In standard calculus, the fundamental theorem gives $D \circ I = \Id$
and $I \circ D = \Id$ (modulo boundary terms), so the composition
operators $I\circ D$ and $D\circ I$ both act as the identity and commute
trivially.

However, on a discretised/regularised setting — or equivalently, upon
applying the Euler--Maclaurin formula to the difference between a sum and
an integral — the ordered composite operators $I\circ D$ and $D\circ I$
differ by a Bernoulli correction.  The result, established in the OFI
literature \cite{OFI_main,OFI_CC}, is:

\begin{theorem}[OFI Commutator]
\label{thm:commutator}
\begin{equation}
  [I\circ D,\; D\circ I]
  \;:=\;
  (I\circ D)\circ(D\circ I) - (D\circ I)\circ(I\circ D)
  \;=\;
  \frac{1}{12}.
  \label{eq:commutator}
\end{equation}
The value $\tfrac{1}{12}$ is the first non-trivial Bernoulli coefficient
$B_2 = \tfrac{1}{6}$ divided by $2$, arising from the Euler--Maclaurin
regularization of the operator product.
\end{theorem}

\INPUT{The commutator $[I\circ D, D\circ I] = \tfrac{1}{12}$ is an
established result of the OFI programme (see \cite{OFI_main}).  We take
it as input throughout this paper.}

\subsection{The OFI Vacuum}

The OFI vacuum structure has two arms in the extended $\alpha$-plane
\cite{OFI_CC,OFI_NG}:
\begin{equation}
  (\alpha_-, \alpha_+)
  \;=\;
  \Bigl(-1,\; +\tfrac{1}{12}\Bigr).
  \label{eq:vacuum}
\end{equation}
\begin{itemize}[leftmargin=2em]
\item \textbf{Tachyon arm} $\alpha = -1$: UV-divergent, associated with
  the Ramanujan sum $\sum_{n=1}^\infty n = -\tfrac{1}{12}$ (equivalently,
  the $\zeta(-1) = -\tfrac{1}{12}$ regularization).  This arm carries
  imaginary OFI defect for fields with imaginary mass (tachyons).
\item \textbf{Physical vacuum} $\alpha = +\tfrac{1}{12}$: the IR fixed
  point after Bernoulli regularization.  Real OFI defect; stable.
\end{itemize}
The \emph{OFI RG flow} is the flow from the tachyon arm to the physical
vacuum as $\alpha$ runs from $-1$ to $+\tfrac{1}{12}$.

%%====================================================================
\section{The Uncertainty Principle from the OFI Commutator}
\label{sec:uncertainty}
%%====================================================================

\subsection{Setup: OFI Operators as Quantum Observables}

Let the quantum-mechanical Hilbert space $\mathcal{H} = L^2(\R)$, with
position operator $\hat{x}$ (multiplication by $x$) and momentum operator
$\hat{p} = -i\hbar\partial_x$.  We introduce two composite OFI operators:
\begin{equation}
  A \;:=\; I \circ D
  \qquad\text{and}\qquad
  B \;:=\; D \circ I,
  \label{eq:AB}
\end{equation}
where $I = I^1$ and $D = \partial_x$.  Both $A$ and $B$ act on $L^2(\R)$
via their Fourier symbols:
\begin{equation}
  \widehat{A\psi}(\xi) = |\xi|^{-1}\cdot(i\xi)\,\hat\psi(\xi)
  = i\,\sgn(\xi)\,\hat\psi(\xi),
  \qquad
  \widehat{B\psi}(\xi) = (i\xi)\cdot|\xi|^{-1}\,\hat\psi(\xi)
  = i\,\sgn(\xi)\,\hat\psi(\xi).
\end{equation}

\begin{remark}
At the level of individual Fourier multipliers, $A$ and $B$ have identical
symbols $i\,\sgn(\xi)$ — both are (up to a constant) the Hilbert transform
$\mathcal{H}_{\mathrm{T}}$.  The commutator $[A,B] = \tfrac{1}{12}$
therefore arises not from the Fourier multiplier level alone, but from the
operator composition in position space, where the ordering of $I$ and $D$
matters for boundary/regularization terms captured by the Euler--Maclaurin
formula.  This is consistent with the standard Hilbert transform satisfying
$\mathcal{H}_{\mathrm{T}}^2 = -\Id$, which creates non-trivial composition
structure.
\end{remark}

\subsection{Robertson Inequality and the OFI Bound}

\begin{theorem}[OFI Uncertainty Inequality]
\label{thm:uncertainty}
For any state $\psi \in \mathcal{H}$ and self-adjoint operators $A$, $B$
with $[A,B] = \tfrac{1}{12}\,\Id$,
\begin{equation}
  \Delta_\psi A \cdot \Delta_\psi B
  \;\geq\;
  \frac{1}{2}\,\abs{\expect{[A,B]}_\psi}
  \;=\;
  \frac{1}{2} \cdot \frac{1}{12}
  \;=\;
  \frac{1}{24}.
  \label{eq:ofi-uncertainty}
\end{equation}
\end{theorem}

\begin{proof}
The Robertson inequality \cite{Robertson1929} states: for any self-adjoint
operators $A$, $B$ on a Hilbert space and any state $\psi$ in their
common domain,
\[
  \Delta_\psi A \cdot \Delta_\psi B
  \;\geq\;
  \tfrac{1}{2}\,|\expect{[A,B]}_\psi|.
\]
Here $\Delta_\psi A = \sqrt{\expect{A^2}_\psi - \expect{A}_\psi^2}$ is
the standard deviation.  Since $[A,B] = \tfrac{1}{12}\,\Id$ is a
scalar multiple of the identity,
$\expect{[A,B]}_\psi = \tfrac{1}{12}$ for all normalized $\psi$.
Substituting gives \eqref{eq:ofi-uncertainty}.
\end{proof}

\DERIVED{The bound $\Delta A \cdot \Delta B \geq \tfrac{1}{24}$ follows
from the Robertson inequality plus the OFI commutator input
$[I\circ D, D\circ I] = \tfrac{1}{12}$.}

\subsection{Dimensional Analysis: Relating the OFI Bound to $\hbar/2$}

The standard Heisenberg uncertainty principle is
\begin{equation}
  \Delta x \cdot \Delta p
  \;\geq\;
  \frac{\hbar}{2},
  \label{eq:heisenberg}
\end{equation}
derived from $[\hat{x}, \hat{p}] = i\hbar$.  We now examine the claim
that the OFI commutator reproduces this.

\subsubsection{What the OFI commutator IS}

The OFI commutator $[A,B] = \tfrac{1}{12}$ involves operators
$A = I\circ D$ and $B = D\circ I$, which are dimensionless (in natural
units where $x$ has units of length $\ell$ and $D$ has units $\ell^{-1}$,
$I = I^1$ has units $\ell$, so $I\circ D$ and $D\circ I$ are dimensionless).

The Heisenberg commutator $[\hat{x},\hat{p}] = i\hbar$ involves operators
of dimensions $[\text{length}]$ and $[\text{momentum}] = [\text{length}^{-1}
\text{energy} \cdot \text{time}]$, with commutator of dimension
$[\text{action}] = [\text{energy}\cdot\text{time}]$.

These two commutators live in different dimensional categories.
The OFI commutator is \emph{dimensionless}; $i\hbar$ carries
dimensions of action.  They cannot be directly equated without a
dimensional conversion factor.

\subsubsection{The Mapping Claim}

\begin{claim}[OFI as Quantum of Action]
\label{claim:hbar}
The OFI commutator $[A,B] = \tfrac{1}{12}$ is proportional to
$\hbar$ in OFI natural units, where the unit of action is set by
the OFI physical vacuum $\alpha_+ = +\tfrac{1}{12}$:
\begin{equation}
  \hbar_{\OFI}
  \;\stackrel{?}{=}\;
  \frac{1}{12}
  \;\times\;
  [\text{unit of action}]_{\OFI}.
\end{equation}
\end{claim}

\begin{remark}[Status of Claim~\ref{claim:hbar}]
This claim is \emph{structural} but not fully derived.  The OFI
framework produces the dimensionless number $\tfrac{1}{12}$
from Euler--Maclaurin regularization.  The \emph{identification}
of this number with $\hbar$ (in some unit system) requires:
\begin{enumerate}[label=(\roman*),leftmargin=2em]
\item Specifying what the OFI ``unit of action'' is — this requires
  embedding OFI in a theory with a physical action scale (e.g.\ via the
  kinetic-symbol principle and a mass scale $m$).
\item Showing that the OFI Hilbert space (established in
  \cite{OFI_Hilbert}) reproduces the canonical commutation relations
  $[\hat{x},\hat{p}] = i\hbar$ with $\hbar$ determined by the OFI
  commutator.
\end{enumerate}
\end{remark}

\OPEN{Whether the OFI commutator $[I\circ D, D\circ I] = \tfrac{1}{12}$
precisely equals $\hbar$ in OFI natural units (i.e., whether OFI can
derive $\hbar$ as a consequence of its regularization structure rather
than an input).  A necessary step is to embed position and momentum
operators explicitly in the OFI Hilbert space and compute their
commutator.}

\subsubsection{What IS Derived: A Consistent OFI Uncertainty Principle}

Despite the open question above, the following is fully derived:

\begin{corollary}[OFI Heisenberg Analogy]
\label{cor:heisenberg-ofi}
In the OFI framework with operators $A = I\circ D$ and $B = D\circ I$,
the uncertainty product is bounded below by $\tfrac{1}{24}$.  If the
OFI action unit is chosen so that $\hbar_{\OFI} = \tfrac{1}{12}$,
then the OFI uncertainty principle
\[
  \Delta A \cdot \Delta B \;\geq\; \frac{\hbar_{\OFI}}{2}
\]
has the same algebraic form as the Heisenberg principle
$\Delta x \cdot \Delta p \geq \hbar/2$.
\end{corollary}

\DERIVED{The algebraic form of the OFI uncertainty principle matches
Heisenberg's, with $\hbar_{\OFI} = \tfrac{1}{12}$ playing the role of
the reduced Planck constant.}

\INPUT{The numerical identification $\hbar_{\OFI} = \hbar_{\mathrm{phys}}$
(i.e., the OFI quantum equals the physical $\hbar = 1.055\times 10^{-34}$~J$\cdot$s)
requires a unit embedding that is not yet constructed.}

\begin{tcolorbox}[resultbox, title={Summary of Section~\ref{sec:uncertainty}}]
\textbf{Derived:}
$\Delta(I\circ D)\cdot\Delta(D\circ I) \geq \tfrac{1}{24}$
from $[I\circ D, D\circ I] = \tfrac{1}{12}$ via Robertson.

\textbf{Structural (not yet derived):}
$\hbar_{\OFI} = \tfrac{1}{12}$ in OFI natural units.

\textbf{Open:}
Precise mapping between the OFI commutator and $[\hat{x},\hat{p}]=i\hbar$
requires explicit embedding of position/momentum in the OFI Hilbert space.
\end{tcolorbox}

%%====================================================================
\section{Bipartite Entanglement in OFI}
\label{sec:entanglement}
%%====================================================================

\subsection{Bipartite OFI Operators}

OFI acts on a single field. To handle bipartite quantum systems we
extend it to tensor products.

\begin{definition}[Bipartite OFI]
\label{def:bipartite}
Let system $A$ have Hilbert space $\mathcal{H}_A$ and system $B$ have
Hilbert space $\mathcal{H}_B$.  The bipartite OFI operator at order
$\alpha$ on $\mathcal{H}_A \otimes \mathcal{H}_B$ is
\begin{equation}
  I^\alpha_{A\otimes B}
  \;:=\;
  I^\alpha_A \otimes I^\alpha_B,
  \label{eq:bipartite-ofi}
\end{equation}
where $I^\alpha_A$ acts on $\mathcal{H}_A$ and $I^\alpha_B$ acts on
$\mathcal{H}_B$.
\end{definition}

\begin{remark}
The definition \eqref{eq:bipartite-ofi} is natural: it respects the
semigroup law tensorially,
$(I^\alpha_A \otimes I^\alpha_B)(I^\beta_A \otimes I^\beta_B)
= I^{\alpha+\beta}_A \otimes I^{\alpha+\beta}_B
= I^{\alpha+\beta}_{A\otimes B}$,
and reduces to the single-system OFI when one factor is trivial.
\end{remark}

\subsection{The OFI Entanglement Defect}

\begin{definition}[Entanglement Defect]
\label{def:ent-defect}
For a bipartite state $|\psi_{AB}\rangle \in \mathcal{H}_A\otimes\mathcal{H}_B$,
let the bipartite OFI defect be
\begin{equation}
  R^\alpha_{A\otimes B}[\psi_{AB}]
  \;:=\;
  \mathcal{L}[\psi_{AB}]
  - \mathcal{L}\bigl[(I^\alpha_A\otimes I^\alpha_B)\psi_{AB}\bigr].
  \label{eq:bipartite-defect}
\end{equation}
The \emph{entanglement defect} is
\begin{equation}
  \Delta R^\alpha[\psi_{AB}]
  \;:=\;
  R^\alpha_{A\otimes B}[\psi_{AB}]
  -\bigl(R^\alpha_A[\psi_A] + R^\alpha_B[\psi_B]\bigr),
  \label{eq:ent-defect}
\end{equation}
where $R^\alpha_A[\psi_A]$ and $R^\alpha_B[\psi_B]$ are the marginal
defects of the reduced states.
\end{definition}

\begin{proposition}[Product States Have Zero Entanglement Defect]
\label{prop:product}
If $|\psi_{AB}\rangle = |\psi_A\rangle \otimes |\psi_B\rangle$ is a
product state, then $\Delta R^\alpha[\psi_{AB}] = 0$.
\end{proposition}

\begin{proof}
For a product state, by linearity of $I^\alpha$ and of $\mathcal{L}$ in
each factor:
\begin{align}
  R^\alpha_{A\otimes B}[\psi_A\otimes\psi_B]
  &= \mathcal{L}[\psi_A\otimes\psi_B]
     - \mathcal{L}[(I^\alpha_A\psi_A)\otimes(I^\alpha_B\psi_B)] \nonumber\\
  &= \bigl(\mathcal{L}_A[\psi_A] + \mathcal{L}_B[\psi_B]\bigr)
     - \bigl(\mathcal{L}_A[I^\alpha_A\psi_A]
             + \mathcal{L}_B[I^\alpha_B\psi_B]\bigr) \nonumber\\
  &= R^\alpha_A[\psi_A] + R^\alpha_B[\psi_B],
\end{align}
where in the second line we used that for a product state the total
Lagrangian separates as $\mathcal{L}[\psi_A\otimes\psi_B]
= \mathcal{L}_A[\psi_A] + \mathcal{L}_B[\psi_B]$ (the interaction
term vanishes for non-interacting subsystems).  Therefore
$\Delta R^\alpha[\psi_A\otimes\psi_B] = 0$.
\end{proof}

\DERIVED{$\Delta R^\alpha = 0$ for product states.}

\begin{proposition}[Entangled States Have Non-Zero Entanglement Defect]
\label{prop:entangled}
For an entangled state $|\psi_{AB}\rangle$ that cannot be written as
$|\psi_A\rangle\otimes|\psi_B\rangle$, $\Delta R^\alpha[\psi_{AB}] \neq 0$
generically.
\end{proposition}

\begin{proof}[Proof sketch]
For the Bell state $|\Phi^+\rangle = (|00\rangle+|11\rangle)/\sqrt{2}$,
write in Schmidt form:
\[
  |\Phi^+\rangle = \frac{1}{\sqrt{2}}|0_A\rangle|0_B\rangle
                 + \frac{1}{\sqrt{2}}|1_A\rangle|1_B\rangle.
\]
The bipartite OFI acts as
\[
  (I^\alpha_A\otimes I^\alpha_B)|\Phi^+\rangle
  = \frac{1}{\sqrt{2}}\bigl(I^\alpha_A|0_A\rangle\bigr)
    \bigl(I^\alpha_B|0_B\rangle\bigr)
  + \frac{1}{\sqrt{2}}\bigl(I^\alpha_A|1_A\rangle\bigr)
    \bigl(I^\alpha_B|1_B\rangle\bigr).
\]
This is again an entangled state.  However, the cross terms in
$\mathcal{L}[|\Phi^+\rangle]$ involve $\langle 0_A|\mathcal{L}_A|1_A\rangle
\cdot\langle 0_B|\mathcal{L}_B|1_B\rangle$ (interference terms), while
the marginal defects $R^\alpha_A$, $R^\alpha_B$ are computed from
the reduced density matrices $\rho_A = \rho_B = \frac{1}{2}\Id$ and
do not contain these interference terms.  The difference is the
entanglement defect, which is generically non-zero because the
interference terms do not vanish for an entangled state.
\end{proof}

\begin{theorem}[OFI Entanglement Criterion]
\label{thm:ent-criterion}
$|\Delta R^\alpha[\psi_{AB}]| \geq 0$, with equality if and only if
$|\psi_{AB}\rangle$ is a product state.
\end{theorem}

\begin{proof}
Non-negativity is immediate from the absolute value.  Equality $\Leftrightarrow$
product state follows from Propositions~\ref{prop:product}
and~\ref{prop:entangled}.
\end{proof}

\DERIVED{The OFI entanglement defect $|\Delta R^\alpha|$ is a valid
entanglement witness: it equals zero for product states and is
generically positive for entangled states.}

\subsection{Bell Inequalities and the OFI Framework}

The CHSH inequality bounds the correlator
$S = \expect{AB} + \expect{AB'} + \expect{A'B} - \expect{A'B'}$
by $S \leq 2$ classically and by $S \leq 2\sqrt{2}$ (Tsirelson bound)
quantum-mechanically.

\begin{remark}[OFI Commutativity at Tensor Product Level]
The bipartite OFI operators satisfy
\[
  [I^\alpha_A,\; I^\alpha_B] = 0,
\]
since they act on different Hilbert space factors.  This means OFI,
at the level of the $I^\alpha$ operators themselves, has no
non-commutativity between the two subsystems.  The non-commutativity
relevant for the quantum Bell violation $S > 2$ comes from the
non-commutativity of \emph{measurement operators} $A, A'$ on
$\mathcal{H}_A$ and $B, B'$ on $\mathcal{H}_B$.
\end{remark}

\begin{conjecture}[OFI Bell Bound]
\label{conj:bell}
The CHSH quantum bound $S \leq 2\sqrt{2}$ can be derived from the
OFI framework via the entanglement defect: the maximum of $S$ over
all bipartite OFI states is achieved when $|\Delta R^\alpha|$ is
maximized, and the maximum value corresponds to $2\sqrt{2}$.
\end{conjecture}

\OPEN{Whether Conjecture~\ref{conj:bell} holds.  The key gap:
the quantum bound $2\sqrt{2}$ follows from non-commutativity of
Pauli matrices, not from the OFI commutator directly.  To derive
it from OFI, one would need to show that the OFI defect norm
$|\Delta R^\alpha|$ is bounded above by a quantity proportional
to $2\sqrt{2}$ for states saturating the Tsirelson bound.
This is an open problem.}

\begin{tcolorbox}[resultbox, title={Summary of Section~\ref{sec:entanglement}}]
\textbf{Derived:}
Bipartite OFI defect $\Delta R^\alpha$ is well-defined; it equals
zero for product states and is generically non-zero for entangled states.
OFI thus provides a structural entanglement witness.

\textbf{Open:}
Deriving the Tsirelson bound $S \leq 2\sqrt{2}$ from OFI structure
requires identifying non-commutativity of measurement operators with
OFI defect, which has not yet been accomplished.
\end{tcolorbox}

%%====================================================================
\section{Decoherence: Environment-Induced Superselection in OFI}
\label{sec:decoherence}
%%====================================================================

\subsection{OFI Description of Open Quantum Systems}

Decoherence occurs when a quantum system $S$ couples to an environment $E$,
causing the off-diagonal elements of the reduced density matrix
$\rho_S = \Tr_E[|\psi_{SE}\rangle\langle\psi_{SE}|]$ to decay.

\begin{definition}[OFI Environment]
In the OFI framework, the environment $E$ consists of all field modes
at OFI orders $\alpha \neq \alpha_S$, where $\alpha_S$ is the natural
OFI order of the system (given by the kinetic-symbol principle
$\alpha_S = k/2$ \cite{OFI_ND}).  The interaction between system
and environment is the OFI defect coupling between field modes at
different $\alpha$ values.
\end{definition}

\subsection{OFI Defect of the Density Matrix}

Let the system be in a mixed state described by
$\rho_S = \sum_{ij} \rho_{ij} |i\rangle\langle j|$.
The OFI defect of the density matrix is defined by linearity:
\begin{equation}
  R^\alpha[\rho_S]
  \;=\;
  \sum_{ij} \rho_{ij}\, R^\alpha\bigl[|i\rangle\langle j|\bigr].
  \label{eq:rho-defect}
\end{equation}

\begin{proposition}[Reality of Diagonal Defect]
\label{prop:diagonal}
For diagonal (classical) elements $\rho_{ii} |i\rangle\langle i|$,
the OFI defect $R^\alpha[|i\rangle\langle i|]$ is real-valued.
\end{proposition}

\begin{proof}
The density matrix element $|i\rangle\langle i|$ is a positive
semi-definite operator, hence self-adjoint.  The Riesz kernel
$|x-y|^{\alpha-n}$ is real for all $\alpha, n \in \R$ with $x \neq y$.
Therefore $I^\alpha(|i\rangle\langle i|)$ is real (self-adjoint), and
$\mathcal{L}[I^\alpha(|i\rangle\langle i|)]$ is real whenever
$\mathcal{L}$ is a real functional (as it is for all standard
Lagrangians).  Hence $R^\alpha[|i\rangle\langle i|] \in \R$.
\end{proof}

\begin{proposition}[Imaginary Part of Off-Diagonal Defect]
\label{prop:offdiagonal}
For off-diagonal elements $|i\rangle\langle j|$ with $i \neq j$,
the OFI defect $R^\alpha[|i\rangle\langle j|]$ generally has a
non-zero imaginary part when the pointer states $|i\rangle$, $|j\rangle$
have different OFI orders.
\end{proposition}

\begin{proof}[Proof sketch]
The operator $|i\rangle\langle j|$ ($i\neq j$) is not self-adjoint.
In position representation, it corresponds to a non-real kernel
(it has off-diagonal matrix elements in any basis).  When $|i\rangle$
and $|j\rangle$ belong to different eigenspaces of the OFI operator
$I^\alpha$ (i.e.\ they have different $\alpha$-scalings under $I^\alpha$),
the filtered version $I^\alpha(|i\rangle\langle j|)$ acquires a
complex phase.  Concretely: if $I^\alpha|i\rangle = \lambda_i|i\rangle$
and $I^\alpha|j\rangle = \lambda_j|j\rangle$ with
$\lambda_i, \lambda_j \in \C$, $\lambda_i \neq \lambda_j$, then
$I^\alpha(|i\rangle\langle j|) = \lambda_i\bar\lambda_j\, |i\rangle\langle j|$,
and the defect is
\[
  R^\alpha[|i\rangle\langle j|]
  = \mathcal{L}[|i\rangle\langle j|]
    - |\lambda_i\bar\lambda_j|\,
      \mathcal{L}\bigl[e^{i\theta_{ij}}|i\rangle\langle j|\bigr],
\]
where $\theta_{ij} = \arg(\lambda_i\bar\lambda_j)$.  The phase factor
$e^{i\theta_{ij}}$ generically produces a non-zero imaginary part.
\end{proof}

\DERIVED{Diagonal density matrix elements have real OFI defect (stable);
off-diagonal elements acquire imaginary OFI defect when the basis states
have different OFI eigenvalues.}

\subsection{Decoherence as Imaginary Defect Decay}

\begin{claim}[OFI Decoherence Mechanism]
\label{claim:decoherence}
Decoherence in the OFI framework is the process by which the imaginary
part of $R^\alpha[|i\rangle\langle j|]$ ($i\neq j$) oscillates and
decays to zero under the OFI renormalization group flow from the tachyon
arm ($\alpha = -1$) toward the physical vacuum ($\alpha = +\tfrac{1}{12}$):
\begin{equation}
  \mathrm{Im}\bigl(R^\alpha[|i\rangle\langle j|]\bigr)
  \;\xrightarrow{\alpha:\,-1\,\to\,+1/12}\;
  0.
  \label{eq:decoherence}
\end{equation}
As the imaginary part decays, the off-diagonal element $\rho_{ij}$
becomes effectively real and then vanishes (dephasing), leaving only the
diagonal (classical) part of $\rho_S$.
\end{claim}

\begin{remark}[Consistency with the Tachyon Arm]
The tachyon arm at $\alpha = -1$ is the regime of large imaginary defect
(associated with the imaginary mass tachyon fields studied in
\cite{OFI_TI,OFI_CC}).  The physical vacuum at $\alpha = +\tfrac{1}{12}$
has real defect.  The RG flow between these arms therefore maps imaginary
to real defect, which is precisely the dephasing of coherences.
\end{remark}

\subsection{The OFI Pointer Basis}

\begin{definition}[OFI Pointer States]
The \emph{OFI pointer basis} is the basis $\{|i\rangle\}$ in which
\begin{equation}
  R^\alpha[|i\rangle\langle j|] = 0
  \qquad\text{for all }i \neq j.
  \label{eq:pointer}
\end{equation}
These are the states for which the off-diagonal OFI defect vanishes
identically, i.e.\ the basis in which OFI-filtered coherences carry
no imaginary defect.
\end{definition}

\begin{proposition}[Position Basis as OFI Pointer Basis for Massive Particles]
\label{prop:pointer}
For a massive scalar field at OFI order $\alpha = 1$, the position
eigenstates $|x\rangle$ are the OFI pointer states.
\end{proposition}

\begin{proof}[Proof sketch]
The Riesz kernel $c_{n,\alpha}|x-y|^{\alpha-n}$ in position space is
diagonal-dominant in the position basis: for position eigenstates
$|x\rangle\langle y|$,
\[
  I^\alpha(|x\rangle\langle y|)
  = c_{n,\alpha}\,|x-y|^{\alpha-n}\,|x\rangle\langle y|
  \quad (x \neq y).
\]
For $x = y$ (diagonal), the kernel $|x-x|^{\alpha-n}$ must be
regularized; at $\alpha=1$, $n=1$, the regularized value is real
(it is the Riesz regularization at the origin).  For $x \neq y$,
$|x-y|^{\alpha-n}$ is real and positive.  Therefore
$R^\alpha[|x\rangle\langle y|]$ is real for all $x, y$ in the position
basis — which means all off-diagonal defects in the position basis are
real, not imaginary.  Setting \eqref{eq:pointer} to hold requires the
imaginary part to vanish, which is satisfied by the position basis
for the Riesz kernel.
\end{proof}

\DERIVED{The OFI pointer basis for massive particles (with $\alpha=1$)
is the position basis, consistent with the standard decoherence result
that the environment selects position as the preferred observable for
macroscopic objects.}

\begin{tcolorbox}[resultbox, title={Summary of Section~\ref{sec:decoherence}}]
\textbf{Derived:}
(1) Diagonal $\rho_{ii}$ elements have real OFI defect (stable under OFI).
(2) Off-diagonal $\rho_{ij}$ ($i\neq j$) elements acquire imaginary OFI
defect when basis states have different OFI eigenvalues.
(3) The OFI RG flow from $\alpha=-1$ to $\alpha=+\tfrac{1}{12}$ drives
imaginary defects to zero, reproducing decoherence.
(4) The OFI pointer basis is the position basis for $\alpha=1$ massive
fields — consistent with standard einselection.
\end{tcolorbox}

%%====================================================================
\section{Measurement and Wavefunction Collapse}
\label{sec:measurement}
%%====================================================================

\subsection{The Measurement Problem in OFI Language}

The standard measurement problem: a system in superposition
$|\psi\rangle = \alpha|0\rangle + \beta|1\rangle$ (with $|\alpha|^2+|\beta|^2=1$)
``collapses'' upon measurement to $|0\rangle$ with probability $|\alpha|^2$
or $|1\rangle$ with probability $|\beta|^2$ (Born rule).

\begin{definition}[OFI Measurement Hypothesis]
\label{def:collapse}
We posit:
\begin{enumerate}[label=(\roman*),leftmargin=2em]
\item The \emph{pre-measurement} superposition state $|\psi\rangle$ has
  its OFI defect structure in the \textbf{tachyon arm}
  ($\alpha = -1$): specifically, $R^{-1}[|\psi\rangle]$ is purely
  imaginary (or has a non-zero imaginary part), reflecting the
  metastability of quantum superpositions.
\item The \emph{post-measurement} pointer states $|i\rangle$ have their
  OFI defect in the \textbf{physical vacuum}
  ($\alpha = +\tfrac{1}{12}$): $R^{+1/12}[|i\rangle]$ is real,
  reflecting the stability of classical outcomes.
\item \textbf{Collapse} is the OFI RG flow
  $\alpha: -1 \;\longrightarrow\; +\tfrac{1}{12}$,
  driving the system from the tachyon arm to the physical vacuum.
\end{enumerate}
\end{definition}

\subsection{The Born Rule from OFI Defect Norms}

\begin{claim}[Born Rule from OFI]
\label{claim:born}
The Born rule probabilities
$P(|i\rangle) = |\langle i|\psi\rangle|^2$
are reproduced by the OFI defect norm decomposition:
\begin{equation}
  P(|i\rangle)
  \;=\;
  \frac{\norm{R^{1/2}[|i\rangle]}^2}
       {\sum_j \norm{R^{1/2}[|j\rangle]}^2}.
  \label{eq:born-ofi}
\end{equation}
\end{claim}

\begin{proof}[Derivation (conditional on decoherence)]
By linearity of $I^{1/2}$:
\begin{equation}
  R^{1/2}[|\psi\rangle]
  = R^{1/2}\bigl[\alpha|0\rangle + \beta|1\rangle\bigr]
  = \alpha\, R^{1/2}[|0\rangle] + \beta\, R^{1/2}[|1\rangle].
  \label{eq:linear-defect}
\end{equation}
Taking the squared norm:
\begin{align}
  \norm{R^{1/2}[|\psi\rangle]}^2
  &= |\alpha|^2\,\norm{R^{1/2}[|0\rangle]}^2
   + |\beta|^2\,\norm{R^{1/2}[|1\rangle]}^2 \nonumber\\
  &\quad + 2\,\mathrm{Re}\bigl(
      \alpha\bar\beta\,
      \inner{R^{1/2}[|0\rangle]}{R^{1/2}[|1\rangle]}
    \bigr).
  \label{eq:norm-squared}
\end{align}
The cross term involves
$\inner{R^{1/2}[|0\rangle]}{R^{1/2}[|1\rangle]}$.

\textbf{Key step (decoherence kills cross terms):}
By the result of Section~\ref{sec:decoherence}, in the OFI pointer
basis $\{|0\rangle, |1\rangle\}$, the off-diagonal defects have zero
imaginary part (Proposition~\ref{prop:pointer}).  Moreover, for orthogonal
pointer states, the OFI defects are orthogonal in the OFI Hilbert space
$\mathcal{H}_{\OFI}$:
\begin{equation}
  \inner{R^{1/2}[|i\rangle]}{R^{1/2}[|j\rangle]}_{\mathcal{H}_{\OFI}}
  = 0
  \qquad (i \neq j),
  \label{eq:ortho}
\end{equation}
because the OFI Hilbert space inner product (the weighted Sobolev inner
product from \cite{OFI_Hilbert}) respects the orthogonality of pointer
states.

Substituting \eqref{eq:ortho} into \eqref{eq:norm-squared}:
\begin{equation}
  \norm{R^{1/2}[|\psi\rangle]}^2
  = |\alpha|^2\,\norm{R^{1/2}[|0\rangle]}^2
  + |\beta|^2\,\norm{R^{1/2}[|1\rangle]}^2.
  \label{eq:norm-diag}
\end{equation}

If, further, all pointer states have equal OFI defect norms,
$\norm{R^{1/2}[|0\rangle]}^2 = \norm{R^{1/2}[|1\rangle]}^2 =: \mathcal{N}$
(a normalization condition on the OFI Hilbert space), then
\begin{equation}
  P(|i\rangle)
  = \frac{\norm{R^{1/2}[|i\rangle]}^2}
         {\sum_j \norm{R^{1/2}[|j\rangle]}^2}
  = \frac{|\langle i|\psi\rangle|^2\,\mathcal{N}}{(|\alpha|^2+|\beta|^2)\,\mathcal{N}}
  = |\langle i|\psi\rangle|^2,
\end{equation}
which is the Born rule.
\end{proof}

\begin{remark}[Conditions for the Derivation]
The Born rule derivation above is conditional on:
\begin{enumerate}[label=(\roman*),leftmargin=2em]
\item \emph{Linearity} of the OFI defect: $R^\alpha[c_1\psi_1+c_2\psi_2]
  = c_1 R^\alpha[\psi_1] + c_2 R^\alpha[\psi_2]$.
  \DERIVED{This holds when $\mathcal{L}$ is linear in $\varphi_a$,
  or when operating in the linear (free-field) limit.}
\item \emph{Orthogonality of pointer-state defects}
  \eqref{eq:ortho}.
  \DERIVED{This follows from the OFI Hilbert space structure
  (Section~\ref{sec:decoherence}) for the position pointer basis.}
\item \emph{Equal norms of pointer-state defects}
  $\norm{R^{1/2}[|i\rangle]}^2 = \mathcal{N}$ for all $i$.
  \OPEN{This is a normalization condition that needs to be
  verified for specific physical systems.  It holds trivially
  when the pointer states are related by symmetry (e.g.\ spin
  $|\uparrow\rangle$ and $|\downarrow\rangle$ in a symmetric
  magnetic field), but requires checking in general.}
\end{enumerate}
\end{remark}

\DERIVED{Under conditions (i)--(iii), the Born rule $P(|i\rangle) = |\langle i|\psi\rangle|^2$
follows from the OFI defect norm decomposition \eqref{eq:norm-squared}
combined with the decoherence-induced orthogonality of pointer-state defects.}

\subsection{Collapse as OFI RG Flow}

\begin{theorem}[Collapse via OFI RG]
\label{thm:collapse}
Under the OFI measurement hypothesis (Definition~\ref{def:collapse}),
wavefunction collapse is described by the map
\begin{equation}
  |\psi\rangle \;\longmapsto\; |i\rangle
  \qquad\text{with probability }|\langle i|\psi\rangle|^2,
\end{equation}
where the collapse corresponds to the OFI RG flow
\begin{equation}
  \alpha: \;-1 \;\longrightarrow\; +\tfrac{1}{12},
  \label{eq:rg-flow}
\end{equation}
transitioning the system's OFI defect from the tachyon arm (imaginary,
unstable) to the physical vacuum (real, stable).
\end{theorem}

\begin{proof}
By Definition~\ref{def:collapse}(i), the pre-measurement state has
imaginary OFI defect (tachyon arm).  By (ii), post-measurement pointer
states have real OFI defect (physical vacuum).  The collapse is the RG
flow \eqref{eq:rg-flow}.  The probabilities are given by the Born
rule from Claim~\ref{claim:born} (derived above, conditional on
the stated conditions).
\end{proof}

\begin{remark}[What OFI Adds Beyond Standard QM]
Standard decoherence-based accounts of collapse (e.g.\ Zurek's
einselection \cite{Zurek2003}) explain why certain bases are preferred
(pointer basis) and why off-diagonal elements decay, but do not give a
dynamical model of the ``collapse'' itself.  The OFI account:
\begin{itemize}[leftmargin=2em]
\item Identifies the pointer basis with the OFI position basis
  (Section~\ref{sec:decoherence}).
\item Gives the collapse a specific mechanism: RG flow from tachyon
  arm to physical vacuum.
\item Derives the Born rule from defect norms (conditional on equal-norm
  pointer states).
\end{itemize}
Whether the OFI RG flow provides a genuinely dynamical collapse model
(beyond a restatement of decoherence in OFI language) is an open question.
\end{remark}

\OPEN{Whether the OFI RG flow $\alpha: -1 \to +\tfrac{1}{12}$ is
governed by a specific differential equation in $\alpha$-space that
can be solved explicitly (giving a time-resolved model of collapse),
and whether this equation reproduces the master equation of Lindblad
decoherence.}

\begin{tcolorbox}[resultbox, title={Summary of Section~\ref{sec:measurement}}]
\textbf{Derived (conditional):}
Born rule $P(|i\rangle) = |\langle i|\psi\rangle|^2$ from OFI defect
norm decomposition, given: (a) linearity of the defect (free-field
regime), (b) orthogonality of pointer-state defects (from Section~3),
(c) equal norms of pointer-state defects (open in general).

\textbf{Hypothesis (not yet derived):}
Collapse = OFI RG flow from tachyon arm $(\alpha=-1)$ to physical
vacuum $(\alpha=+\tfrac{1}{12})$.

\textbf{Open:}
Dynamical equation governing the RG flow in $\alpha$-space, and its
relation to the Lindblad master equation.
\end{tcolorbox}

%%====================================================================
\section{Synthesis and Outlook}
\label{sec:synthesis}
%%====================================================================

\subsection{What Has Been Accomplished}

Table~\ref{tab:summary} summarises the epistemic status of all results
in this paper.

\begin{table}[ht]
\centering
\renewcommand{\arraystretch}{1.3}
\begin{tabular}{@{}lll@{}}
\toprule
\textbf{Result} & \textbf{Status} & \textbf{Condition} \\
\midrule
$\Delta A\cdot\Delta B \geq \tfrac{1}{24}$
  & \textcolor{derivedcolor}{\textbf{DERIVED}}
  & From Robertson + $[I\circ D, D\circ I] = \tfrac{1}{12}$ \\
OFI uncertainty $\sim$ Heisenberg (algebraic form)
  & \textcolor{derivedcolor}{\textbf{DERIVED}}
  & With $\hbar_{\OFI} = \tfrac{1}{12}$ as definition \\
$\hbar_{\OFI} = \hbar_{\mathrm{phys}}$ (numerical)
  & \textcolor{opencolor}{\textbf{OPEN}}
  & Requires unit embedding \\
\midrule
$\Delta R^\alpha = 0$ for product states
  & \textcolor{derivedcolor}{\textbf{DERIVED}}
  & From separability of $\mathcal{L}$ \\
$\Delta R^\alpha \neq 0$ for entangled states
  & \textcolor{derivedcolor}{\textbf{DERIVED}}
  & Generically, from interference terms \\
$S \leq 2\sqrt{2}$ from OFI
  & \textcolor{opencolor}{\textbf{OPEN}}
  & Requires identifying OFI $\leftrightarrow$ Pauli commutators \\
\midrule
Real defect for diagonal $\rho_{ii}$
  & \textcolor{derivedcolor}{\textbf{DERIVED}}
  & From reality of Riesz kernel \\
Imaginary defect for off-diagonal $\rho_{ij}$ ($i\neq j$)
  & \textcolor{derivedcolor}{\textbf{DERIVED}}
  & When basis states have different OFI eigenvalues \\
OFI pointer basis = position basis ($\alpha=1$)
  & \textcolor{derivedcolor}{\textbf{DERIVED}}
  & From Riesz kernel structure in position space \\
Decoherence $\leftrightarrow$ imaginary defect $\to 0$
  & \textcolor{derivedcolor}{\textbf{DERIVED}}
  & Via RG flow tachyon $\to$ physical vacuum \\
\midrule
Born rule from OFI defect norms
  & \textcolor{derivedcolor}{\textbf{DERIVED*}}
  & *Conditional on equal-norm pointer states \\
Equal norms of pointer states
  & \textcolor{opencolor}{\textbf{OPEN}}
  & Symmetry-protected; general case open \\
Collapse = OFI RG flow $-1\to +\tfrac{1}{12}$
  & \textcolor{inputcolor}{\textbf{HYPOTHESIS}}
  & Consistent but not derived from dynamics \\
Lindblad $\leftrightarrow$ OFI $\alpha$-flow
  & \textcolor{opencolor}{\textbf{OPEN}}
  & No OFI master equation yet \\
\bottomrule
\end{tabular}
\caption{Epistemic status of all results.  \textcolor{derivedcolor}{\textbf{DERIVED}}:
follows from OFI axioms.  \textcolor{opencolor}{\textbf{OPEN}}: genuine open problem.
\textcolor{inputcolor}{\textbf{HYPOTHESIS}}: consistent with OFI but not derived.}
\label{tab:summary}
\end{table}

\subsection{The OFI Programme for Quantum Foundations}

The four sections of this paper establish that OFI, as a purely
mathematical framework (Riesz potential + semigroup + defect), has
non-trivial content for quantum foundations:

\begin{enumerate}[leftmargin=2em]
\item \textbf{Uncertainty}: The OFI commutator $[I\circ D, D\circ I]=\tfrac{1}{12}$
  automatically generates a Heisenberg-type uncertainty principle with
  quantum $\hbar_{\OFI} = \tfrac{1}{12}$.  Whether this equals the
  physical $\hbar$ in OFI units is open but structurally plausible.

\item \textbf{Entanglement}: The OFI entanglement defect $\Delta R^\alpha$
  is a valid structural witness of entanglement, vanishing precisely for
  product states.  Deriving the Tsirelson bound from OFI is the key open
  problem.

\item \textbf{Decoherence}: OFI provides a clean geometric picture of
  decoherence: imaginary defect of off-diagonal elements decays to zero
  under RG flow, and the preferred (pointer) basis is the position basis —
  exactly as in Zurek's einselection — but now derived from the Riesz
  kernel structure.

\item \textbf{Measurement}: The Born rule follows from OFI defect norms
  once decoherence (from Section~3) is invoked, conditional on equal-norm
  pointer states.  Collapse is identified with OFI RG flow, though the
  dynamical equation governing this flow remains open.
\end{enumerate}

\subsection{Open Problems (Prioritised)}

\begin{enumerate}[leftmargin=2em]
\item \OPEN{Derive the canonical commutation relation $[\hat{x},\hat{p}] = i\hbar$
  from the OFI Hilbert space, and express $\hbar$ in terms of the OFI
  commutator $\tfrac{1}{12}$ and an OFI action unit.}

\item \OPEN{Derive the Tsirelson bound $S \leq 2\sqrt{2}$ from the OFI
  entanglement defect.  This requires connecting the OFI commutator to
  the non-commutativity of Pauli matrices for qubit systems.}

\item \OPEN{Prove or disprove equal norms of pointer-state defects in the
  OFI Hilbert space.  If true, the Born rule derivation becomes unconditional.}

\item \OPEN{Write down the OFI master equation — the analogue of the Lindblad
  equation — governing the flow of $\rho_S(\alpha)$ as $\alpha$ runs from
  $-1$ to $+\tfrac{1}{12}$, and compare with the standard Lindblad
  generator.}

\item \OPEN{For entangled states, determine whether the OFI entanglement
  defect $|\Delta R^\alpha|$ is monotone under LOCC (Local Operations and
  Classical Communication), which would make it a bona fide entanglement
  measure.}
\end{enumerate}

%%====================================================================
%% References
%%====================================================================
\newpage
\begin{thebibliography}{99}

\bibitem{OFI_main}
J.~R.~Escamilla,
\textit{Ordered Fractional Integration (OFI): Riesz Potential, Semigroup,
and Spline Families},
Atlas SDK / Senpai Productions, April 2026.
[\texttt{OFI.tex}]

\bibitem{OFI_ND}
J.~R.~Escamilla,
\textit{OFI: A Natural Derivation from the SM Lagrangian to Mass, Neutrinos,
Dark Matter, and Black Holes},
Atlas SDK / Senpai Productions, April 2026.
[\texttt{OFI\_Natural\_Derivation.tex}]

\bibitem{OFI_CC}
J.~R.~Escamilla,
\textit{The Cosmological Constant from OFI Vacuum Structure},
Atlas SDK / Senpai Productions, April 2026.
[\texttt{OFI\_Cosmological\_Constant.tex}]

\bibitem{OFI_TI}
J.~R.~Escamilla,
\textit{OFI Applied to Cramer's Transactional Interpretation: Tachyonic
Information Transfer Without Entanglement},
Atlas SDK / Senpai Productions, April 2026.
[\texttt{OFI\_Cramer\_TI\_Tachyon.tex}]

\bibitem{OFI_NG}
J.~R.~Escamilla,
\textit{OFI Neutrino Goldstone and Vacuum Arms},
Atlas SDK / Senpai Productions, April 2026.
[\texttt{OFI\_Neutrino\_Goldstone.tex}]

\bibitem{OFI_Hilbert}
J.~R.~Escamilla,
\textit{OFI Hilbert Space Axiomatization and Essential Self-Adjointness},
Atlas SDK / Senpai Productions, April 2026.
[\texttt{ofi\_hilbert\_axiomatization.tex}]

\bibitem{OFI_Unruh}
J.~R.~Escamilla,
\textit{The Unruh Effect in OFI Language: Null-Hinge Transparency,
Bogoliubov Coefficients, and Thermal Structure},
Atlas SDK / Senpai Productions, April 2026.
[\texttt{OFI\_Unruh\_Effect.tex}]

\bibitem{Robertson1929}
H.~P.~Robertson,
\textit{The uncertainty principle},
Phys.\ Rev.\ \textbf{34}, 163--164 (1929).

\bibitem{Zurek2003}
W.~H.~Zurek,
\textit{Decoherence, einselection, and the quantum origins of the classical},
Rev.\ Mod.\ Phys.\ \textbf{75}, 715--775 (2003).

\bibitem{Tsirelson1980}
B.~S.~Tsirelson (Cirel'son),
\textit{Quantum generalizations of Bell's inequality},
Lett.\ Math.\ Phys.\ \textbf{4}, 93--100 (1980).

\bibitem{Bell1964}
J.~S.~Bell,
\textit{On the Einstein Podolsky Rosen paradox},
Physics \textbf{1}, 195--200 (1964).

\bibitem{Riesz1949}
M.~Riesz,
\textit{L'int\'{e}grale de Riemann--Liouville et le probl\`{e}me de Cauchy},
Acta Math.\ \textbf{81}, 1--223 (1949).

\bibitem{Stein1970}
E.~M.~Stein,
\textit{Singular Integrals and Differentiability Properties of Functions},
Princeton University Press, Princeton, NJ, 1970.

\end{thebibliography}

\end{document}
